

% 3. Environnement unique et superposition directe
\begin{tikzpicture}
  % On dessine les deux éléments sur le MÊME repère absolu
  \clip[](-4, -4) rectangle (11 , 4) ;

  \node[transform canvas={scale=0.5}] at (15,0) {
    \begin{tikzpicture}

        \draw[
            opacity = 0.5, 
            fill=colorOne!10!colorTwo,
            draw=colorOne,
            line width=2.2ex
            ] (0,0) ellipse [x radius=5cm, y radius=2cm];

        \draw[
            fill=none,
            draw=colorOne,
            line width=2.2ex
            ] (0,1) ellipse [x radius=5cm, y radius=2cm];

        \draw[
            color=colorOne,
            line width=2.2ex,
            pattern={Lines[
                        angle=90,
                        distance=3ex,
                        line width=1.8ex
                    ]},
            pattern color=colorOne!50!colorTwo ,
            ] (5,0) arc[
            start angle=0,
            end angle=-180,
            x radius=5cm,
            y radius=2cm
            ]--(-5,1)
            (-5,1) arc[
            start angle=180,
            end angle=360,
            x radius=5cm,
            y radius=2cm
            ]--(5,0)
            (5,0) arc[
            start angle=0,
            end angle=180,
            x radius=5cm,
            y radius=2cm
            ]--(-5,1)
            (-5,1) arc[
            start angle=180,
            end angle=0,
            x radius=5cm,
            y radius=2cm
            ]--(5,0)
        ;

        \draw[
            color=colorFour,
            line width=1.0ex,
            ->
            ] (0,0) -- ({5*cos(45)},{2*sin(45)})
            node[pos = 0.5, above , font=\fontsize{20pt}{1pt}\selectfont\bfseries, text=colorFour  , rotate = 30 ]{$r$}
        ;

        \draw[
            fill=colorFour,
            color=colorFour,
            line width=0.5ex,
            ] (0,0) circle (0.1)
        ;

         \draw[
            color=colorFour,
            line width=0.5ex,
            >-
            ] (-6,0) -- (-6,0.2)
        ;
        \draw[
            color=colorFour,
            line width=0.5ex,
            -<
            ] (-6,0.8) -- (-6,1)
        ;
        \node[font=\fontsize{10pt}{1pt}\selectfont\bfseries, text=colorFour  , rotate = 90 ] at (-6, 0.5) {$\mathrm{d} r$}; 

        %\node[font=\fontsize{40pt}{1pt}\selectfont\bfseries, text=colorFour  , rotate = 0 ] at (0, -3) {$\mathrm{d} r \not \ll r $};

        \begin{scope}[shift={(-4,4)}]
                \draw[,
                    line width=1.5ex, color=colorOne, 
                fill=colorSix!20,
                    pattern={Lines[
                        angle=90,
                        distance=3ex,
                        line width=1.8ex
                    ]},
                    pattern color=colorOne!50!colorTwo ]  
                (0,0) rectangle (3,2);
                \node[font=\fontsize{40pt}{1pt}\selectfont\bfseries, text=colorOne , right ] 
                    at (3.5,1) {$=$};

                \node[font=\fontsize{40pt}{1pt}\selectfont\bfseries, text=colorOne , right ] 
                    at (5.5,1) {$\mathrm{d}A(r)$};


        \end{scope}

    \end{tikzpicture}   
  };

  \node[transform canvas={scale=0.5}] at (0,0) {
    \begin{tikzpicture}

        \begin{scope}

            \draw[
            color=colorSix,
            line width=2.2ex,
            draw = none, 
            pattern={Lines[
                        angle=45,
                        distance=3ex,
                        line width=1.8ex
                    ]},
            pattern color=colorSix ,
            ] (5,0) -- (6,0) 
            (6,0) arc[
            start angle=0,
            end angle=360,
            x radius=6cm,
            y radius=6cm
            ]--(5,0)
            (5,0) arc[
            start angle=0,
            end angle=-360,
            x radius=5cm,
            y radius=5cm
            ]
        ;

        \end{scope}

        \draw[
            fill=none,
            draw=colorOne,
            line width=2.2ex
            ] (0,0) circle (5);

        \draw[
            fill=none,
            draw=colorOne,
            line width=2.2ex
            ] (0,0) circle (6);
        
       

        \draw[
            color=colorFour,
            line width=1.0ex,
            ->
            ] (0,0) -- ({5*cos(45)},{5*sin(45)})
            node[pos = 0.5, above , font=\fontsize{20pt}{1pt}\selectfont\bfseries, text=colorFour  , rotate = 45 ]{$r$}
        ;

        \draw[
            fill=colorFour,
            color=colorFour,
            line width=0.5ex,
            ] (0,0) circle (0.1)
        ;

         \draw[
            color=colorFour,
            line width=0.5ex,
            >-
            ] (-6,0) -- (-5.8,0.)
        ;
        \draw[
            color=colorFour,
            line width=0.5ex,
            -<
            ] (-5.2,0) -- (-5,0)
        ;
        \node[font=\fontsize{10pt}{1pt}\selectfont\bfseries, text=colorFour   , rotate = 0 ] at (-5.5, 0.0) {$\mathrm{d} r$}; 

        %\node[font=\fontsize{40pt}{1pt}\selectfont\bfseries, text=colorFour  , rotate = 0 ] at (0, -7) {$\mathrm{d} r  \ll r $};

        \begin{scope}[shift={(-4,6.5)}]
                \draw[,
                    line width=1.5ex, color=colorOne, 
                fill=colorSix!20,
                    pattern={Lines[
                        angle=45,
                        distance=3ex,
                        line width=1.8ex
                    ]},
                    pattern color=colorSix ]  
                (0,0) rectangle (3,2);
                \node[font=\fontsize{40pt}{1pt}\selectfont\bfseries, text=colorOne , right ] 
                    at (3.5,1) {$=$};

                \node[font=\fontsize{40pt}{1pt}\selectfont\bfseries, text=colorOne , right ] 
                    at (5.5,1) {$\Delta A(r)$};


        \end{scope}

    \end{tikzpicture}   
  };


  

  % Grille de repérage
  %\draw[line width=2.5ex , color = red] (-8, -5.5) rectangle (8 , 4.5) ;
  %\drawgrid{-3}{11}{-4}{4}{1}{1}
\end{tikzpicture}